% Amplificador com transistor PNP da Questão 2 (Figura fig:q2-circuito).
% Compilar com pdflatex; o PDF resultante substitui figuras/fig_q2_circuito.pdf.
% Entra no relatório em tamanho natural (escala 1).
\documentclass[11pt,border=4pt]{standalone}
% Estilo comum dos esquemáticos do relatório (circuitikz).
% Uso: \documentclass[11pt,border=4pt]{standalone} e, em seguida, % Estilo comum dos esquemáticos do relatório (circuitikz).
% Uso: \documentclass[11pt,border=4pt]{standalone} e, em seguida, % Estilo comum dos esquemáticos do relatório (circuitikz).
% Uso: \documentclass[11pt,border=4pt]{standalone} e, em seguida, \input{estilo_circuitos}.
\usepackage[T1]{fontenc}
\usepackage[utf8]{inputenc}
\usepackage[brazilian]{babel}
\usepackage{lmodern}
\usepackage{amsmath}
\usepackage[american]{circuitikz}
\usetikzlibrary{calc,arrows.meta}

% Cores das figuras originais
\definecolor{anot}{RGB}{0,135,200}      % anotações: valores calculados, setas, resistências vistas
\definecolor{sombra}{RGB}{220,242,253}  % região do transistor nos modelos T e pi
\definecolor{term}{RGB}{62,124,148}     % nomes de terminais destacados (B, C, E)

% Fios com 0,4 pt e corpo dos componentes com 0,8 pt, como nas figuras originais
\ctikzset{bipoles/thickness=2}
.
\usepackage[T1]{fontenc}
\usepackage[utf8]{inputenc}
\usepackage[brazilian]{babel}
\usepackage{lmodern}
\usepackage{amsmath}
\usepackage[american]{circuitikz}
\usetikzlibrary{calc,arrows.meta}

% Cores das figuras originais
\definecolor{anot}{RGB}{0,135,200}      % anotações: valores calculados, setas, resistências vistas
\definecolor{sombra}{RGB}{220,242,253}  % região do transistor nos modelos T e pi
\definecolor{term}{RGB}{62,124,148}     % nomes de terminais destacados (B, C, E)

% Fios com 0,4 pt e corpo dos componentes com 0,8 pt, como nas figuras originais
\ctikzset{bipoles/thickness=2}
.
\usepackage[T1]{fontenc}
\usepackage[utf8]{inputenc}
\usepackage[brazilian]{babel}
\usepackage{lmodern}
\usepackage{amsmath}
\usepackage[american]{circuitikz}
\usetikzlibrary{calc,arrows.meta}

% Cores das figuras originais
\definecolor{anot}{RGB}{0,135,200}      % anotações: valores calculados, setas, resistências vistas
\definecolor{sombra}{RGB}{220,242,253}  % região do transistor nos modelos T e pi
\definecolor{term}{RGB}{62,124,148}     % nomes de terminais destacados (B, C, E)

% Fios com 0,4 pt e corpo dos componentes com 0,8 pt, como nas figuras originais
\ctikzset{bipoles/thickness=2}

\begin{document}
\begin{circuitikz}
  % Folga entre rótulo e traço e meias-alturas dos símbolos no tamanho padrão do circuitikz
  \def\g{6pt}        % folga desejada
  \def\zz{5.98pt}    % meia-altura do zigue-zague do resistor
  \def\pl{11.95pt}   % meia-altura das placas do capacitor e da placa longa da bateria; raio da fonte

  % Trilho superior: v_i, R_S, C_C1, nó de R_E e emissor de Q
  \draw (0,-3.1) node[ground]{} to[sV] (0,0)
        to[R] (2.8,0)
        to[C] (5.0,0) node[circ]{}
        -- (6.1,0) node[pnp, rotate=90, anchor=E] (Q) {};
  % Coletor de Q, nó de R_C, C_C2, nó de R_L e saída
  \draw (Q.C) -- (8.34,0) node[circ]{}
        to[C] (11.34,0) node[circ]{}
        -- (12.34,0) node[ocirc]{};

  % Base de Q no terra e ramos verticais
  % (o traço longo da bateria fica do lado do início do caminho)
  \draw (Q.B) node[ground]{};
  \draw (5.0,0) to[R] (5.0,-2.4) to[battery1] (5.0,-3.8) node[ground]{};
  \draw (8.34,-3.8) node[ground]{} to[battery1] (8.34,-2.4) to[R] (8.34,0);
  \draw (11.34,0) to[R] (11.34,-3.8) node[ground]{};

  % Rótulos
  \node[anchor=east, inner sep=0] at ($(0,-1.55)+(-\pl-\g,0)$) {$v_i$};
  \node[anchor=south, inner sep=0] at ($(1.4,0)+(0,\zz+\g+1.3pt)$) {$R_S$};
  \node[anchor=north, inner sep=0] at ($(1.4,0)+(0,-\zz-\g)$) {$1\,\mathrm{k\Omega}$};
  \node[anchor=south, inner sep=0] at ($(3.9,0)+(0,\pl+\g+1.5pt)$) {$C_{C1}$};
  \node[anchor=south, inner sep=0] at ($($(Q.E)!0.5!(Q.C)$)+(0,\pl+\g+3.5pt)$) {$Q$};
  \node[anchor=east, align=left, inner sep=0] at ($(5.0,-1.2)+(-\zz-\g,0)$) {$R_E$\\$10\,\mathrm{k\Omega}$};
  \node[anchor=east, inner sep=0] at ($(5.0,-3.1)+(-\pl-\g,0)$) {$20\,\mathrm{V}$};
  \node[anchor=west, align=left, inner sep=0] at ($(8.34,-1.2)+(\zz+\g,0)$) {$R_C$\\$6{,}5\,\mathrm{k\Omega}$};
  \node[anchor=west, inner sep=0] at ($(8.34,-3.1)+(\pl+\g,0)$) {$25\,\mathrm{V}$};
  \node[anchor=south, inner sep=0] at ($(9.84,0)+(0,\pl+\g+1.5pt)$) {$C_{C2}$};
  \node[anchor=west, align=left, inner sep=0] at ($(11.34,-1.9)+(\zz+\g,0)$) {$R_L$\\$5\,\mathrm{k\Omega}$};
  \node[anchor=west, inner sep=0] at ($(12.34,0)+(1.6pt+\g,0)$) {$v_o$};
\end{circuitikz}
\end{document}
